\documentclass[conference,twocolumn,spanish]{IEEEtran}

\usepackage[utf8]{inputenc}
\usepackage[T1]{fontenc}
\usepackage[spanish,es-tabla]{babel}
\usepackage{graphicx}
\usepackage{fancyhdr}
\usepackage{booktabs}
\usepackage{tabularx}
\usepackage{pgfplots}
\pgfplotsset{compat=1.18}
\usepackage{hyperref}
\usepackage{cite}

% Referencia ficticia autocontenida para visualizar cómo se renderiza la cita en el texto
\begin{filecontents*}{referencias_plantilla.bib}
@misc{ref_ejemplo_2026,
  author       = {{Centro de Inteligencia Artificial e Innovación (CIAI)}},
  title        = {Documento técnico de referencia de ejemplo para plantilla institucional},
  year         = {2026},
  howpublished = {\url{https://www.uts.edu.co}}
}
\end{filecontents*}

% Metadatos internos del documento PDF (incluye derechos y licencia CC BY-NC-ND 4.0)
\hypersetup{
    pdftitle={Boletín | Centro de Inteligencia Artificial e Innovación -- Plantilla},
    pdfauthor={Miguel Ángel Osorio Orduz; Jhoan Sebastian Torres Garcia; Frank Camilo Zuñiga Vargas; Docentes: Carlos Alfonso Pineda Ortiz, Diego Alexander Cruz Melo},
    pdfsubject={Boletín Semanal de Vigilancia Tecnológica en Inteligencia Artificial (CIAI - UTS)},
    pdfkeywords={Inteligencia Artificial, LLM, CIAI, UTS, Vigilancia Tecnológica, Benchmarks},
    pdfcreator={LaTeX con clase IEEEtran},
    pdfproducer={Centro de Inteligencia Artificial e Innovación (CIAI) - Unidades Tecnológicas de Santander},
    pdfcopyright={Copyright (C) 2026 CIAI - UTS. Licencia Creative Commons Atribución-NoComercial-SinDerivadas 4.0 Internacional (CC BY-NC-ND 4.0).},
    pdflicenseurl={https://creativecommons.org/licenses/by-nc-nd/4.0/deed.es}
}

% Configuración de encabezado y pie de página
\pagestyle{fancy}
\fancyhf{}
\fancyhead[L]{\small Boletín Semanal CIAI -- Plantilla}
\fancyhead[R]{\includegraphics[height=0.8cm]{assets/uts_logonegro.png}}
\fancyfoot[L]{\scriptsize \copyright~2026 CIAI -- UTS | Licencia \href{https://creativecommons.org/licenses/by-nc-nd/4.0/deed.es}{CC BY-NC-ND 4.0} (Uso académico, no comercial, sin derivadas)}
\fancyfoot[R]{\scriptsize \thepage}
\renewcommand{\headrulewidth}{0.4pt}
\renewcommand{\footrulewidth}{0.2pt}

% Macro dinámica para gráficas de avance (0% a 100%) de los Proyectos CIAI
% Uso: \graficaavanceproyecto{Título corto}{Hito4,Hito3,Hito2,Hito1}{{(Hito4,0) (Hito3,25) (Hito2,70) (Hito1,100)}}
\newcommand{\graficaavanceproyecto}[3]{%
\begin{figure}[htbp]
\centering
\begin{tikzpicture}
\begin{axis}[
    xbar,
    width=0.92\columnwidth,
    height=3.6cm,
    bar width=7pt,
    xlabel={Porcentaje de avance (\%)},
    xmin=0,
    xmax=114,
    xtick={0,25,50,75,100},
    symbolic y coords={#2},
    ytick=data,
    tick label style={font=\scriptsize},
    label style={font=\scriptsize},
    xmajorgrids=true,
    grid style={dashed, gray!35},
    nodes near coords={\pgfmathprintnumber{\pgfplotspointmeta}\%},
    every node near coord/.append style={font=\tiny, anchor=west}
]
\addplot[fill=green!55!black, draw=green!70!black] coordinates {#3};
\end{axis}
\end{tikzpicture}
\caption{Porcentaje de avance por hitos: #1.}
\end{figure}
}

\title{Boletín | Centro de Inteligencia Artificial e Innovación\\
       \large Plantilla}

\author{
  \IEEEauthorblockN{Miguel Ángel Osorio Orduz}
  \IEEEauthorblockN{Jhoan Sebastian Torres Garcia}
  \IEEEauthorblockN{Frank Camilo Zuñiga Vargas}
  \IEEEauthorblockA{Docentes: Carlos Alfonso Pineda Ortiz, Diego Alexander Cruz Melo}
  \IEEEauthorblockA{
    Centro de Inteligencia Artificial e Innovación (CIAI)\\
    Colombia\\
    mangelosorio@uts.edu.co\\
    jhoanstorres@uts.edu.co\\
    fczuniga@uts.edu.co
  }
}

\begin{document}

\maketitle
\thispagestyle{fancy}

\begin{abstract}
Nota de ejemplo para el resumen ejecutivo de la plantilla del boletín semanal del CIAI~\cite{ref_ejemplo_2026}.
\end{abstract}

% =========================================================
\section{Investigación}
% =========================================================
Dato de ejemplo plantilla~\cite{ref_ejemplo_2026}. A continuación se ilustran los estilos visuales de gráficas disponibles (barras, histograma y circular) con valores genéricos de demostración.

\begin{figure}[htbp]
\centering
\begin{tikzpicture}
\begin{axis}[
    ybar,
    width=\columnwidth,
    height=4.4cm,
    bar width=11pt,
    ylabel={Eje Y (Unidad)},
    symbolic x coords={Cat A, Cat B, Cat C, Cat D},
    xtick=data,
    tick label style={font=\scriptsize},
    label style={font=\scriptsize},
    ymajorgrids=true,
    grid style={dashed, gray!35},
    ymin=0,
    ymax=100,
    yticklabels={}
]
\addplot[fill=blue!65!black, draw=blue!80!black] coordinates {(Cat A,75) (Cat B,50) (Cat C,85) (Cat D,40)};
\end{axis}
\end{tikzpicture}
\caption{Ejemplo visual de gráfica de barras verticales en doble columna.}
\end{figure}

\begin{figure}[htbp]
\centering
\begin{tikzpicture}
\begin{axis}[
    ybar interval,
    width=\columnwidth,
    height=4.4cm,
    ylabel={Frecuencia},
    xlabel={Intervalos de clase},
    xmin=0,
    xmax=5,
    ymin=0,
    ymax=100,
    xtick={0,1,2,3,4,5},
    xticklabels={I1, I2, I3, I4, I5, },
    yticklabels={},
    tick label style={font=\scriptsize},
    label style={font=\scriptsize},
    ymajorgrids=true,
    grid style={dashed, gray!35}
]
\addplot[fill=teal!65!black, draw=teal!85!black] coordinates {(0,30) (1,65) (2,90) (3,55) (4,25) (5,0)};
\end{axis}
\end{tikzpicture}
\caption{Ejemplo visual de histograma de distribución sin valores numéricos.}
\end{figure}

\begin{figure}[htbp]
\centering
\begin{tikzpicture}[scale=0.85]
% Gráfico circular nativo en TikZ (sin paquetes externos)
\def\radius{1.45cm}
\fill[blue!70!black] (0,0) -- (0:\radius) arc (0:144:\radius) -- cycle;
\fill[teal!70!black] (0,0) -- (144:\radius) arc (144:252:\radius) -- cycle;
\fill[orange!80!black] (0,0) -- (252:\radius) arc (252:324:\radius) -- cycle;
\fill[gray!65] (0,0) -- (324:\radius) arc (324:360:\radius) -- cycle;
\draw[white, thick] (0,0) -- (0:\radius) (0,0) -- (144:\radius) (0,0) -- (252:\radius) (0,0) -- (324:\radius);
\draw[black!40] (0,0) circle (\radius);

% Leyenda lateral
\node[anchor=west, font=\scriptsize] at (1.8, 0.75) {\textcolor{blue!70!black}{$\blacksquare$} Segmento A};
\node[anchor=west, font=\scriptsize] at (1.8, 0.25) {\textcolor{teal!70!black}{$\blacksquare$} Segmento B};
\node[anchor=west, font=\scriptsize] at (1.8, -0.25) {\textcolor{orange!80!black}{$\blacksquare$} Segmento C};
\node[anchor=west, font=\scriptsize] at (1.8, -0.75) {\textcolor{gray!65}{$\blacksquare$} Segmento D};
\end{tikzpicture}
\caption{Ejemplo visual de gráfica circular (torta) nativa en TikZ.}
\end{figure}

% =========================================================
\section{Herramientas}
% =========================================================
Dato de ejemplo plantilla~\cite{ref_ejemplo_2026}.

% =========================================================
\section{Proyectos CIAI}
% =========================================================
Dato de ejemplo plantilla~\cite{ref_ejemplo_2026}.

% Inserción dinámica de avances de proyectos internos del CIAI mediante archivos externos
% Ejemplo de uso: \input{proyectos/nombre_proyecto}

% =========================================================
\section{Noticias}
% =========================================================
Dato de ejemplo plantilla~\cite{ref_ejemplo_2026}.

\bibliographystyle{IEEEtran}
\bibliography{referencias}

\end{document}